\clearpage
\setcounter{secnumdepth}{-1}
\section{\hfil Abstract \hfil}
\vspace{0.5cm}
\thispagestyle{front}

Hallucination remains a major limitation of retrieval-augmented generation (RAG), since retrieved evidence does not guarantee that a language model will generate a faithful response. Existing hallucination mitigation methods often operate as \textbf{open-loop} detectors or external verifiers: they estimate hallucination risk, but do not directly connect that estimate to inference-time response control. This work proposes a \textbf{closed-loop RAG framework} that links internal-state hallucination detection to automated generation control. To our knowledge, this is the first closed-loop RAG framework that connects dual internal activation probes from frozen SwiGLU decoder LLMs to a three-action inference-time controller for hallucination mitigation. The system instruments frozen decoder-only language models with PyTorch forward hooks and extracts two complementary activation signals: Contextualized Embedding Vectors (CEV) from the residual stream and Intermediate Activation Vectors (IAV) from the SwiGLU feed-forward sublayer before the down-projection layer. Separate lightweight MLP probes are trained on RAGTruth hallucination labels and fused into a calibrated hallucination probability. This probability drives a three-action controller: ACCEPT, INTERVENE, or ABSTAIN; where intervention triggers corrective generation behavior such as regeneration or retrieval revision before a final response is delivered. The framework is evaluated on Mistral-7B-Instruct-v0.2, Qwen3-8B, and LLaMA-3-8B-Instruct using Colab A100 FP16 experiments. On a stratified RAGTruth validation split, the fused probe achieves AUROC scores of \textbf{0.8596} for Mistral-7B-Instruct-v0.2, \textbf{0.8808} for Qwen3-8B, and \textbf{0.8689} for LLaMA-3-8B-Instruct, surpassing all published and reproduced baseline detectors considered in this study. Closed-loop intervention reduces the delivered hallucination proxy relative to vanilla RAG for all three backbones, which span a clear safety–utility spectrum: Mistral-7B-Instruct-v0.2 reaches the lowest delivered proxy at a balanced acceptance rate, Qwen3-8B gives the largest relative reduction but is the most conservative, and LLaMA-3-8B-Instruct retains the most answers. These results suggest that internal activation signals can support practical inference-time hallucination control across SwiGLU-based decoder families, while probe training and controller thresholds remain backbone-specific due to model-dependent activation distributions.
\\

\textbf{Keywords:} Retrieval-Augmented Generation, Hallucination Mitigation, Internal Activation Probing, Closed-Loop Inference Control, Transformer Representations, RAGTruth.

\setcounter{secnumdepth}{1}
\clearpage